% ============================================================================
%  অধ্যায় ০ — শুরুর আগে
%  COMPILE WITH XeLaTeX  (twice, so the point total resolves)
%  Overleaf: Menu -> Compiler -> XeLaTeX
% ============================================================================
\documentclass[11pt]{article}
\usepackage{bnhandout}          % add [nosol] for a student edition

\hauthor[https://jugantordey.github.io/]{Jugantor Dey}
\hseries{\textit{Mathematics for the Physics Olympiad}}

\begin{document}

\handouttitle{Chapter 0 : Put on Your Life Jacket!}

\noindent
এই অধ্যায়ের কোনো পূর্বশর্ত নেই। পঞ্চম শ্রেণির পাটিগণিত জানা থাকলেই যথেষ্ট :
গুণ ও ভাগ, গুণনীয়ক ও গুণিতক, ল.সা.গু ও গ.সা.গু, সাধারণ ভগ্নাংশ ও দশমিক ভগ্নাংশ,
শতকরা, গড়, এবং বন্ধনীসহ যোগ-বিয়োগ-গুণ-ভাগ হিসাবের ক্রম। এগুলো এখানে আর দেখানো হবে না; তুমি
এগুলো পারো ধরে নেওয়া হচ্ছে।

\noindent
তাহলে এই অধ্যায়ে কী আছে? ঠিক সেই কয়টা জিনিস যা পঞ্চম শ্রেণিতে শেখানো
\emph{হয় না}, অথচ এই সিরিজের দ্বিতীয় অধ্যায় থেকেই বারবার লাগবে। ঋণাত্মক
সংখ্যা, বর্গমূলের চিহ্ন, মূলদ ও অমূলদের পার্থক্য, আর ``অসংজ্ঞায়িত'' কথাটার
মানে। চারটাই ছোট বিষয়, কিন্তু একটাও না জানলে পরের অধ্যায়গুলোতে হঠাৎ
আটকে যেতে হয়। এখানে মোট \totalpoints\ পয়েন্ট আছে।

\section{ঋণাত্মক সংখ্যা ও সংখ্যারেখা}

\begin{idea}[সংখ্যারেখা দুদিকেই চলে]
এতদিন সংখ্যারেখা মানে ছিল $0$ থেকে শুরু করে ডান দিকে $1, 2, 3, \dots$। কিন্তু
রেখাটা বাঁ দিকে না যাওয়ার তো কোনো কারণ নেই। $0$-এর বাঁয়ে সমান দূরত্বে যে বিন্দুগুলো
পড়ে, তাদের নাম $-1, -2, -3, \dots$।

এখানে ``ঋণাত্মক'' মানে সংখ্যাটা খারাপ  কিছু নয়। এটা একটা \emph{দিক}।
$5$ মানে $0$ থেকে ডানে পাঁচ ধাপ, $-5$ মানে $0$ থেকে বাঁয়ে পাঁচ ধাপ। দূরত্ব
একই, দিক উল্টো।
\end{idea}

\noindent
দিকের কথাটা মনে রাখলে ক্রম নিয়ে আর বিভ্রান্তি হয় না। সংখ্যারেখায় যে যত ডানে,
সে তত বড়। তাই $-5 < -2$, যদিও $5 > 2$। পাঁচ ধাপ বাঁয়ে যাওয়া দুই ধাপ বাঁয়ে
যাওয়ার চেয়ে বেশি বাঁয়ে, তাই -2 এর চেয়ে -5 ছোট।

\begin{example}
$-7$, $3$, $0$, $-1$, $-10$, $2$ সংখ্যাগুলোকে ছোট থেকে বড় সাজাও।
\end{example}

\begin{solbox}
সংখ্যারেখায় বসিয়ে বাঁ থেকে ডানে পড়লেই হলো :
\[
  -10 \;<\; -7 \;<\; -1 \;<\; 0 \;<\; 2 \;<\; 3 .
\]
লক্ষ করো, ঋণাত্মকগুলোর মধ্যে যেটার সংখ্যামান বড় ($10$ বনাম $7$), সেটাই ছোট।
\end{solbox}

\begin{idea}[বিয়োগ মানে বিপরীত সংখ্যা যোগ]
প্রতিটি সংখ্যা $a$-এর (ধরে নাও, $a$ একটি সংখ্যা) একটি বিপরীত সংখ্যা আছে, লেখা হয় $-a$, যার সংজ্ঞা হলো
$a + (-a) = 0$। সংখ্যারেখায় $a$ আর $-a$ শূন্যের দুই পাশে সমান দূরত্বে বসে।

এতে বিয়োগ আলাদা কোনো কাজ থাকে না :
\[
  a - b \;=\; a + (-b) .
\]
$7 - 10$ মানে সংখ্যারেখায় ডানে সাত ধাপ গিয়ে তারপর বাঁয়ে দশ ধাপ, অর্থাৎ $-3$-এ পৌঁছানো।
আগে বলা হতো ``বড় থেকে ছোট বিয়োগ করতে হয়''; সেই নিষেধটা এখন উঠে গেল।
\end{idea}

\begin{remark}
$-a$ লেখাটা দেখে ভেবে বসার দরকার নেই যে সংখ্যাটা ঋণাত্মক। $a = -4$ হলে
$-a = 4$, যা ধনাত্মক। $-a$ মানে ``$a$-এর বিপরীত'', ``একটি ঋণাত্মক সংখ্যা'' নয়।
এই ভুলটা বীজগণিতে খুব মারাত্মক।
\end{remark}

\noindent
এবার গুণ। স্কুলে সাধারণত তিনটা নিয়ম মুখস্থ করানো হয়, তাদের মধ্যে ``ঋণাত্মক গুণ ঋণাত্মক
সমান ধনাত্মক'' কথাটা শুনতে খামখেয়ালি লাগে। কিন্তু এটা আসলে কেউ ইচ্ছামতো বানায়নি; বরং এটা ছাড়া অন্য কিছু
হলে পাটিগণিত ভেঙে পড়ত। নিচের তালিকাটা দেখো।

\begin{idea}[চিহ্নের নিয়ম আবিষ্কার নয়, বাধ্যবাধকতা]
$-1$ দিয়ে গুণ করার তালিকাটা উপর থেকে নিচে লিখি :
\[
  3 \times (-1) = -3 \quad
\]
\[
  2 \times (-1) = -2 \quad
\]
\[
  1 \times (-1) = -1  \quad
\]
\[
  0 \times (-1) = 0 \quad
\]
বাঁ পাশের সংখ্যা প্রতিবার $1$ করে কমছে, আর ডান পাশের ফল প্রতিবার $1$ করে
\emph{বাড়ছে}। এর পরের ধাপে নিশ্চয়ই -1 দিয়ে গুণ করব, তখন  বাঁ পাশে বসবে $-1$, তাই ডান পাশে বসতে হবে $1$ :
\[
  (-1) \times (-1) = 1 .
\]
এটাই একমাত্র মান যাতে ধারাটা ভাঙে না। একইভাবে সাধারণ নিয়ম দাঁড়ায় : দুটি
সংখ্যার চিহ্ন এক হলে গুণফল ধনাত্মক, চিহ্ন আলাদা হলে ঋণাত্মক। ভাগের বেলাতেও
একই কথা, কারণ ভাগ গুণেরই উল্টো কাজ।
\end{idea}

\begin{example}
হিসাব করো : (ক) $(-6) \times (-4)$, \; (খ) $(-6) + (-4)$, \;
(গ) $(-20) \div 5$, \; (ঘ) $3 - (-8)$।
\end{example}

\begin{solbox}
(ক) চিহ্ন এক, তাই গুণফল ধনাত্মক : $24$।

(খ) এটা গুণ নয়, যোগ। বাঁয়ে ছয় ধাপ, তারপর আরও বাঁয়ে চার ধাপ : $-10$।
চিহ্নের নিয়ম এখানে খাটে না; ওটা গুণ ও ভাগের নিয়ম।

(গ) চিহ্ন আলাদা, তাই ফল ঋণাত্মক : $-4$।

(ঘ) $3 - (-8) = 3 + 8 = 11$। বিপরীত সংখ্যার বিপরীত হলো সংখ্যাটাই।
\end{solbox}

\section{বর্গমূল}

\begin{idea}[বর্গমূল একটি প্রশ্ন]
$\sqrt{9}$ চিহ্নটা কোনো হিসাবের নির্দেশ নয়, একটা প্রশ্ন : ``কোন সংখ্যাকে বর্গ
করলে $9$ পাওয়া যায়?''

উত্তর হতে পারত $3$ বা $-3$, দুটোরই বর্গ $9$। কিন্তু একটা চিহ্নের দুটো মান হলে
হযবরল হওয়ার চান্স আছে। তাই রীতি হলো : $\sqrt{\;}$ চিহ্ন সবসময় \emph{অঋণাত্মক}
উত্তরটা বোঝায়। অর্থাৎ $\sqrt{9} = 3$, কখনোই $-3$ নয়।

অন্য মানটা দরকার হলে সামনে চিহ্ন বসিয়ে লিখতে হয় : $-\sqrt{9} = -3$, আর দুটোই
একসাথে বোঝাতে $\pm\sqrt{9}$।
\end{idea}

\noindent
কিছু সংখ্যার বর্গমূল সহজেই আসে, কারণ তারা পূর্ণবর্গ :
\[
  \sqrt{1}=1, \; \sqrt{4}=2, \; \sqrt{9}=3, \; \sqrt{16}=4, \; \sqrt{25}=5, \; \sqrt{36}=6, \; \sqrt{49}=7, \;
  \]\[ \sqrt{64}=8, \; \sqrt{81}=9, \; \sqrt{100}=10, \; 
\sqrt{121}=11, \; \sqrt{144}=12 
\]
এই তালিকাটা মুখস্থ থাকলে পরের অধ্যায়গুলোতে অনেক সময় বাঁচবে।

\begin{idea}[পূর্ণবর্গ না হলেও বর্গমূল থাকে]
$2$ পূর্ণবর্গ নয়, তবু $\sqrt{2}$ একটা সংখ্যা। কোন সংখ্যা, সেটা আন্দাজ করা যায় :
$1^2 = 1$ আর $2^2 = 4$, অর্থাৎ যে সংখ্যার বর্গ $2$, সেটা $1$ আর $2$-এর মাঝে
পড়বে। আরও কাছে গেলে $1.4^2 = 1.96$ এবং $1.5^2 = 2.25$, তাই $\sqrt{2}$ আছে
$1.4$ ও $1.5$-এর মাঝখানে।

এভাবে যত খুশি নিখুঁত করা যায়, কিন্তু কখনোই ঠিক শেষ হয় না। কেন শেষ হয় না,
সেটা পরের অনুচ্ছেদের বিষয়।
\end{idea}

\begin{remark}
অঋণাত্মক প্রতিটি সংখ্যার এমন একটি বর্গমূল সত্যিই \emph{আছে} কি না, সেই
প্রশ্নটা এখানে তোলা হচ্ছে না; ধরে নেওয়া হচ্ছে। অধ্যায় ৩-এ দ্বিঘাত
সমীকরণের সূত্র বের করার সময় এই ধরে নেওয়াটা আবার স্পষ্ট করে বলা হবে। আপাতত
শুধু একটা কথা মনে রাখো : ঋণাত্মক সংখ্যার বর্গমূল কোনো বাস্তব সংখ্যা নয়, কারণ
যেকোনো বাস্তব সংখ্যার বর্গই অঋণাত্মক (হয় 0 না হয় ধনাত্মক)। $\sqrt{-4}$ লিখলে তুমি বাস্তব সংখ্যার
জগৎ ছেড়ে বেরিয়ে যাচ্ছ; সেই জগতের দেখা পাবে অধ্যায় ১৮-তে।
\end{remark}

\section{মূলদ ও অমূলদ সংখ্যা}

\begin{idea}[মূলদ সংখ্যা মানে ভগ্নাংশে লেখা যায়]
যে সংখ্যাকে $\dfrac{p}{q}$ আকারে লেখা যায়, যেখানে $p$ ও $q$ পূর্ণসংখ্যা এবং
$q \neq 0$, তাকে বলে মূলদ সংখ্যা (rational number)।

তালিকাটা যত ভাবো তার চেয়ে বড়। $\dfrac{3}{4}$ তো বটেই; কিন্তু $7 = \dfrac{7}{1}$,
$-2 = \dfrac{-2}{1}$, $0 = \dfrac{0}{1}$, এমনকি $0.25 = \dfrac{1}{4}$ ও
$0.333\ldots = \dfrac{1}{3}$ সবই মূলদ। অর্থাৎ সব পূর্ণসংখ্যা মূলদ, সব সাধারণ
ভগ্নাংশ মূলদ।
\end{idea}

\begin{idea}[চেনার সহজ উপায় : দশমিক রূপ]
কোনো সংখ্যা মূলদ কি না, তার দশমিক রূপ দেখলেই বোঝা যায়। মূলদ সংখ্যার দশমিক রূপ
দুরকমের একটা হবেই:

\begin{enumerate}
  \item \textbf{শেষ হয়ে যায় :} $\dfrac{1}{4} = 0.25$, \; $\dfrac{7}{8} = 0.875$।
  \item \textbf{পুনরাবৃত্তি হয় :} $\dfrac{1}{3} = 0.333\ldots$, \;
        $\dfrac{2}{11} = 0.181818\ldots$ । একই গুচ্ছ অঙ্ক বারবার ফিরে আসে।
\end{enumerate}

\noindent
এর বাইরে তৃতীয় কোনো সম্ভাবনা মূলদ সংখ্যার নেই। $p$-কে $q$ দ্বারা ভাগ করতে গেলে ভাগশেষ
$0$ থেকে $q-1$ এর মধ্যে কয়েকটা মান হতে  পারে, তাই হয় ভাগশেষ একসময়
$0$ হয়ে যাবে, নয়তো আগের কোনো ভাগশেষ ফিরে আসবে এবং তখন থেকে সব পুনরাবৃত্তি
হতে থাকবে।
\end{idea}

\begin{idea}[অমূলদ সংখ্যা]
যে সংখ্যার দশমিক রূপ শেষও হয় না, পুনরাবৃত্তিও হয় না, তাকে বলে অমূলদ সংখ্যা
(irrational number)। এদের কোনো $\dfrac{p}{q}$ রূপ নেই।

সবচেয়ে চেনা দুটি উদাহরণ :
\[
  \sqrt{2} = 1.41421356\ldots, \qquad \pi = 3.14159265\ldots
\]
দুটোরই অঙ্ক অসীম, এবং কোথাও কোনো পুনরাবৃত্ত গুচ্ছ নেই।

মূলদ ও অমূলদ মিলিয়ে যা পাওয়া যায়, তাদের বলা হয় বাস্তব সংখ্যা (real number)।
সংখ্যারেখার প্রতিটি বিন্দু একটি বাস্তব সংখ্যা, এবং প্রতিটি বাস্তব সংখ্যা
রেখার একটি বিন্দু।
\end{idea}

\begin{remark}
$\sqrt{2}$ ও $\pi$ যে সত্যিই অমূলদ, তার প্রমাণ এখানে দেওয়া হচ্ছে না; ফলটা
মেনে নেওয়া হচ্ছে। প্রমাণ দুটো এই সিরিজের কাজে লাগে না, তাই বাদ রইল। তবে ফলটার
একটা ব্যবহারিক দিক আছে : $\pi$-এর বদলে $22/7$ বা $3.1416$ লিখলে সবসময়ই কিছুটা
ভুল থেকে যাচ্ছে, কারণ কোনো ভগ্নাংশ দিয়েই একে হুবহু লেখা সম্ভব নয়। তাই সত্যিই যদি $\pi$-এর মান দিয়ে নিখুঁত হিসাব করতে চাও, তাহলে অনন্তকাল ধরে তোমাকে শুধু $\pi$-এর মানই লিখতে হবে!\newline
সংখ্যাগুলোকে set-এর ভাষায় $\mathbb{N}$, $\mathbb{Z}$, $\mathbb{Q}$,
$\mathbb{R}$ চিহ্ন দিয়ে গুছিয়ে লেখা হয়। সেটা আসবে অধ্যায় ৪-এ, function-এর
প্রয়োজনে। এখানে শুধু সংখ্যাগুলো চিনে রাখাই যথেষ্ট।
\end{remark}

\begin{example}
নিচের সংখ্যাগুলো মূলদ না অমূলদ, কারণসহ বলো :
(ক) $0.75$, \; (খ) $\sqrt{16}$, \; (গ) $\sqrt{5}$, \; (ঘ) $0.121212\ldots$
\end{example}

\begin{solbox}
(ক) মূলদ। দশমিক শেষ হয়ে গেছে; $0.75 = \dfrac{3}{4}$।

(খ) মূলদ। $16$ পূর্ণবর্গ, তাই $\sqrt{16} = 4 = \dfrac{4}{1}$। বর্গমূলের চিহ্ন
দেখলেই অমূলদ ভাবা ভুল।

(গ) অমূলদ। $5$ পূর্ণবর্গ নয়; $2^2 = 4$ ও $3^2 = 9$, তাই $\sqrt{5}$ পড়ছে $2$
আর $3$-এর মাঝে, কোনো পূর্ণসংখ্যা নয়। পূর্ণসংখ্যার বর্গমূল হয় পূর্ণসংখ্যা,
নয় অমূলদ; মাঝামাঝি কিছু হয় না।

(ঘ) মূলদ। $12$ গুচ্ছটা পুনরাবৃত্তি হচ্ছে; আসলে $\dfrac{4}{33}$।
\end{solbox}

\section{অসংজ্ঞায়িত, অনির্ণেয় ও অসীম}

\noindent
এই অনুচ্ছেদটা ছোট, কিন্তু এখানকার ভুলটা সবচেয়ে বেশি প্রচলিত। এমনকি ইউনিভার্সিটির বহু শিক্ষার্থীও এই প্রশ্নের ভুল উত্তর দিয়ে থাকে! প্রশ্নটা হলো :
$\dfrac{1}{0}$ কত?

\begin{idea}[ভাগের আসল সংজ্ঞা]
$\dfrac{a}{b}$ মানে কী? মানে হলো : এমন একটি সংখ্যা $c$, যার জন্য
$c \times b = a$।

$\dfrac{6}{2} = 3$, কারণ $3 \times 2 = 6$। ভাগ আলাদা কোনো কাজ নয়, গুণের
উল্টো কাজই ভাগ। এই সংজ্ঞাটা হাতে থাকলে শূন্য দিয়ে ভাগের প্রশ্নটার উত্তর নিজে
নিজেই বেরিয়ে আসে।
\end{idea}

\begin{idea}[$1/0$ কেন অসংজ্ঞায়িত]
$\dfrac{1}{0}$ মানে এমন একটি সংখ্যা $c$ খোঁজা, যার জন্য $c \times 0 = 1$।

কিন্তু যেকোনো সংখ্যাকে $0$ দিয়ে গুণ করলে $0$ পাওয়া যায়। তাই $c \times 0$
কখনোই $1$ হবে না, $c$ যা-ই হোক। এমন কোনো সংখ্যা \emph{নেই}।

সংখ্যাটা কঠিন নয়, বড় নয়, লুকানো নয়। এজগতে এরকম কোনো সংখ্যাই নেই যে এই প্রশ্নের উত্তর হতে পারে। তাই $\dfrac{1}{0}$-কে বলা হয়
অসংজ্ঞায়িত (undefined)।
\end{idea}

\begin{idea}[$1/0 = \infty$ কথাটা ভুল]
অনেকেই লেখে $\dfrac{1}{0} = \infty$। এটা ভুল, এবং ভুলটা দুই ধাপের।

প্রথমত, $\infty$ কোনো সংখ্যা নয়। এটা একটা চিহ্ন, একটা ধারণা, যা দিয়ে বোঝানো হয় ``কোনো
সীমা ছাড়িয়ে বাড়তে থাকা''। সংখ্যারেখায় $\infty$ নামের কোনো বিন্দু নেই।
$\infty$-কে সংখ্যা ধরে হিসাব করতে গেলে সঙ্গে সঙ্গে অসংগতি বেরিয়ে আসে।

দ্বিতীয়ত, ভুলটা এসেছে একটা সত্য পর্যবেক্ষণ থেকে :
\[
  \frac{1}{0.1} = 10, \qquad
  \frac{1}{0.01} = 100, \qquad
  \frac{1}{0.001} = 1000, \qquad \dots
\]
হর যত $0$-এর কাছে যাচ্ছে, ভাগফল তত বড় হচ্ছে। কথাটা ঠিক। কিন্তু ``কাছে যাচ্ছে''
আর ``সমান'' এক জিনিস নয়। উপরের তালিকায় হর কখনোই $0$ হয়নি।

আর হর যদি সংখ্যারেখার বাঁ দিক থেকে $0$-এর কাছে যায়, অর্থাৎ $-0.1, -0.01, -0.001$, তাহলে
ভাগফল হয় $-10, -100, -1000$, অর্থাৎ ছোট হতে হতে নামছে। একই দিকে যাচ্ছে না।
তাই ``$1/0$-এর মান অসীম'' কথাটা কোন দিক থেকে আসছি তার উপর নির্ভর করে, এবং
সেজন্যই একে কোনো একটা মান দেওয়া যায় না।
\end{idea}

\begin{remark}
$\dfrac{0}{0}$ আলাদা রকমের ব্যর্থতা। এখানে $c \times 0 = 0$ সমীকরণটা
\emph{সব} $c$-এর জন্যই সত্য। আগেরটায় কোনো উত্তর ছিল না, এখানে সব সংখ্যাই
উত্তর। একটা মান বেছে নেওয়ার কোনো যুক্তি নেই, তাই এটাও অসংজ্ঞায়িত।

এই দ্বিতীয় ধরনের ব্যর্থতার জন্য গণিতে একটা আলাদা নাম আছে : অনির্ণেয়
(indeterminate)। কিন্তু নামটার আসল মানে বোঝার জন্য limit লাগে, যা আসবে
অধ্যায় ১৪-তে। আপাতত এটুকুই যথেষ্ট : পাটিগণিতে $\dfrac{1}{0}$ আর
$\dfrac{0}{0}$ দুটোই অসংজ্ঞায়িত, কোনোটারই মান নেই, এবং কোনোটাই $\infty$ নয়।
\end{remark}

\section{কোণের পরিমাপ : তিনটি মাইলফলক}

\noindent
কোণ মাপা হয় ডিগ্রিতে, আর মাপার যন্ত্র চাঁদা (protractor)। এটুকু আগের থেকেই জানো আশা করি। শুধু তিনটা সংখ্যা এখানে স্পষ্ট করে বলে রাখি, কারণ ত্রিকোণমিতিতে
(অধ্যায় ৭) এদের বারবার লাগবে।

\begin{idea}[$90^\circ$, $180^\circ$, $360^\circ$]
একটা পূর্ণ ঘূর্ণন, অর্থাৎ ঘুরে ঠিক আগের জায়গায় ফিরে আসা, হলো $360^\circ$।
বাকি দুটো এর অংশ :

\begin{enumerate}
  \item পূর্ণ ঘূর্ণনের অর্ধেক $180^\circ$। এতে দিক ঠিক উল্টে যায়, আর দুই বাহু
        মিলে একটা সরলরেখা তৈরি করে। একে বলে সরলকোণ।
  \item পূর্ণ ঘূর্ণনের এক-চতুর্থাংশ $90^\circ$। একে বলে সমকোণ, আর তখন বাহু
        দুটোকে বলে পরস্পর লম্ব।
\end{enumerate}

\noindent
$360$ সংখ্যাটা প্রকৃতির কোনো নিয়ম নয়, নিছক একটা পুরোনো রীতি; সুবিধা এই যে
$360$-এর গুণনীয়ক অনেক, তাই ভাগ করলে সহজে পূর্ণসংখ্যা পাওয়া যায়। অধ্যায়
৭-এ দেখবে ডিগ্রির বদলে radian নামের আরেকটা একক আছে, যেটা ক্যালকুলাসের জন্য
ডিগ্রির চেয়ে অনেক ভালো।
\end{idea}

\section{সমস্যা}

\prob{1} নিচের সংখ্যাগুলোকে সংখ্যারেখায় বসাও এবং ছোট থেকে বড় সাজাও :
$\;-3$, \; $1$, \; $-8$, \; $0$, \; $-1$, \; $5$, \; $-12$।

\begin{sol}
সংখ্যারেখায় বাঁ থেকে ডানে পড়লে :
\[
  -12 \;<\; -8 \;<\; -3 \;<\; -1 \;<\; 0 \;<\; 1 \;<\; 5 .
\]
\end{sol}

\prob{2} হিসাব করো :
(ক) $(-3) \times (-7)$, \;
(খ) $(-3) + (-7)$, \;
(গ) $(-3) - (-7)$, \;
(ঘ) $(-42) \div (-6)$, \;
(ঙ) $(-2) \times 3 \times (-4)$।

\begin{sol}
(ক) চিহ্ন এক, গুণফল ধনাত্মক : $21$।

(খ) যোগ, চিহ্নের নিয়ম খাটে না : $-10$।

(গ) $(-3) + 7 = 4$।

(ঘ) চিহ্ন এক : $7$।

(ঙ) দুবারে করি। $(-2) \times 3 = -6$, তারপর $(-6) \times (-4) = 24$।
ঋণাত্মক সংখ্যা জোড়সংখ্যক হলে গুণফল ধনাত্মক, বিজোড়সংখ্যক হলে ঋণাত্মক।
\end{sol}

\prob{2} ``$-a$ সবসময় একটি ঋণাত্মক সংখ্যা'' \; এই দাবিটা কি সত্য? সত্য হলে
কারণ দেখাও, মিথ্যা হলে একটি counterexample দাও।

\begin{sol}
মিথ্যা। $a = -5$ ধরি। তাহলে $-a = -(-5) = 5$, যা ধনাত্মক। একটিমাত্র
counterexample-ই দাবিটা ভাঙার জন্য যথেষ্ট।

আসল কথা : $-a$ মানে ``$a$-এর বিপরীত''। $a$ ঋণাত্মক হলে তার বিপরীত ধনাত্মক।
আর $a = 0$ হলে $-a = 0$, যা ধনাত্মকও নয় ঋণাত্মকও নয়।
\end{sol}

\prob{1} $\sqrt{144}$, \; $\sqrt{81}$, \; $-\sqrt{25}$ \; এবং $\pm\sqrt{49}$
নির্ণয় করো।

\begin{sol}
$\sqrt{144} = 12$, \; $\sqrt{81} = 9$, \; $-\sqrt{25} = -5$, \;
$\pm\sqrt{49} = 7$ অথবা $-7$।

লক্ষ করো $\sqrt{144}$-এর উত্তর $\pm 12$ নয়। চিহ্নটা একাই অঋণাত্মক মূলটা বোঝায়;
অন্যটা চাইলে সামনে $-$ বা $\pm$ বসাতে হয়।
\end{sol}

\prob{2} $\sqrt{30}$ কোন দুটি পরপর পূর্ণসংখ্যার মাঝে পড়ে? এরপর দশমিকের এক ঘর
পর্যন্ত বলো এটি কোন দুটি সংখ্যার মাঝে আছে।

\begin{sol}
$5^2 = 25$ এবং $6^2 = 36$, আর $25 < 30 < 36$। তাই
\[
  5 \;<\; \sqrt{30} \;<\; 6 .
\]
এবার কাছাকাছি যাই : $5.4^2 = 29.16$ এবং $5.5^2 = 30.25$। যেহেতু
$29.16 < 30 < 30.25$, তাই
\[
  5.4 \;<\; \sqrt{30} \;<\; 5.5 .
\]
\end{sol}

\prob{2} নিচের প্রতিটি সংখ্যা মূলদ না অমূলদ, কারণসহ বলো :
(ক) $\sqrt{64}$, \; (খ) $\sqrt{10}$, \; (গ) $-\dfrac{7}{3}$, \;
(ঘ) $2.5$, \; (ঙ) $0.454545\ldots$

\begin{sol}
(ক) মূলদ। $\sqrt{64} = 8$, একটি পূর্ণসংখ্যা।

(খ) অমূলদ। $10$ পূর্ণবর্গ নয়, তাই এর বর্গমূল কোনো পূর্ণসংখ্যা নয়।

(গ) মূলদ। ইতিমধ্যেই $p/q$ আকারে আছে; ঋণাত্মক হলেও মূলদ।

(ঘ) মূলদ। দশমিক শেষ হয়ে গেছে; $2.5 = \dfrac{5}{2}$।

(ঙ) মূলদ। $45$ গুচ্ছটা পুনরাবৃত্ত।
\end{sol}

\prob{2} একজন শিক্ষার্থী বলল : ``দশমিকের পর অসীম অঙ্ক থাকলেই সংখ্যাটা অমূলদ।''
সে ঠিক বলেছে কি? একটি counterexample দিয়ে উত্তর দাও।

\begin{sol}
ভুল বলেছে। $\dfrac{1}{3} = 0.3333\ldots$ -এর দশমিক অঙ্ক অসীম, তবু এটি মূলদ,
কারণ একে $p/q$ আকারে লেখা গেছে।

আসল শর্তটা দুই ভাগের : অমূলদ হতে হলে দশমিক অসীম \emph{এবং} অপুনরাবৃত্ত হতে
হবে। শুধু অসীম হওয়া যথেষ্ট নয়।
\end{sol}

\prob{3} ভাগের সংজ্ঞা ব্যবহার করে দেখাও যে $\dfrac{5}{0}$-এর কোনো মান হতে পারে
না। এরপর দেখাও যে $\dfrac{0}{0}$-এর ব্যর্থতাটা আলাদা রকমের।

\begin{sol}
\textbf{$5/0$।} সংজ্ঞা অনুযায়ী $\dfrac{5}{0}$ মানে এমন সংখ্যা $c$, যার জন্য
$c \times 0 = 5$। কিন্তু যেকোনো সংখ্যার সাথে $0$ গুণ করলে $0$ হয়, তাই
$c \times 0 = 5$ কোনো $c$-এর জন্যই সত্য নয়। উত্তর হতে পারে এমন সংখ্যার
সংখ্যা শূন্য।

\medskip
\textbf{$0/0$।} এখানে খুঁজতে হবে এমন $c$, যার জন্য $c \times 0 = 0$। এটা
\emph{সব} $c$-এর জন্যই সত্য : $c = 1$ চলে, $c = 7$ চলে, $c = -103$ চলেও।
উত্তর হতে পারে এমন সংখ্যা অসীম সংখ্যক।

\medskip
দুটোই অসংজ্ঞায়িত, কিন্তু কারণ উল্টো : প্রথমটায় কোনো উত্তর নেই, দ্বিতীয়টায়
এত উত্তর যে কোনো একটাকে বেছে নেওয়ার যুক্তি নেই।
\end{sol}

\prob{2} নিচের যুক্তিটিতে ভুল কোথায় খুঁজে বের করো।

\smallskip
\noindent
ধরি $a = b$। তাহলে $a^2 = ab$, তাই $a^2 - b^2 = ab - b^2$, অর্থাৎ
$(a+b)(a-b) = b(a-b)$। দুই পাশকে $(a-b)$ দিয়ে ভাগ করে পাই $a + b = b$।
$a = b = 1$ বসিয়ে পাই $2 = 1$।

\begin{sol}
ভুলটা ভাগের ধাপে। $a = b$ ধরা হয়েছে, তাই $a - b = 0$। অর্থাৎ দুই পাশকে
$(a-b)$ দিয়ে ভাগ করা মানে $0$ দিয়ে ভাগ করা, যা অসংজ্ঞায়িত। ধাপটা বৈধ নয়,
তাই তার পরের কোনো কিছুই মানে না।

এই ফাঁদটা এত সাধারণ যে মনে রাখার মতো : কোনো রাশি দিয়ে ভাগ করার আগে সবসময়
নিশ্চিত হও যে রাশিটা শূন্য নয়।
\end{sol}

\prob{2} একজন লিখল $\dfrac{1}{0} = \infty$, তারপর দুই পাশে $0$ দিয়ে গুণ করে
লিখল $1 = \infty \times 0 = 0$। কোথায় গোলমাল হলো?

\begin{sol}
গোলমাল প্রথম লাইনেই। $\dfrac{1}{0}$-এর কোনো মান নেই, তাই ওটাকে $\infty$-এর
সমান লেখা যায় না; আর $\infty$ কোনো সংখ্যা নয় বলে তাকে নিয়ে গুণ-ভাগও করা যায় না।

একবার একটা অর্থহীন লেখাকে সমীকরণে বসিয়ে ফেললে তার পরে যেকোনো কিছু ``প্রমাণ''
করে ফেলা যায়, যেমন এখানে $1 = 0$ বেরিয়ে এসেছে। ভুল ফলটাই ইঙ্গিত দিচ্ছে যে
শুরুর ধাপটা বৈধ ছিল না।
\end{sol}

\prob{2} একটি ঘড়ির মিনিটের কাঁটা $2$টা থেকে $2$টা $15$ মিনিট পর্যন্ত ঘুরল।
কাঁটাটি কত ডিগ্রি ঘুরল? $2$টা থেকে $2$টা $45$ মিনিট পর্যন্ত ঘুরলে কত হতো?

\begin{sol}
পূর্ণ এক ঘণ্টায় মিনিটের কাঁটা এক পাক, অর্থাৎ $360^\circ$ ঘোরে।

$15$ মিনিট এক ঘণ্টার এক-চতুর্থাংশ, তাই ঘূর্ণন
$\dfrac{1}{4} \times 360^\circ = 90^\circ$, একটি সমকোণ।

$45$ মিনিট এক ঘণ্টার তিন-চতুর্থাংশ, তাই ঘূর্ণন
$\dfrac{3}{4} \times 360^\circ = 270^\circ$।
\end{sol}

\prob{3} $-4$ থেকে শুরু করে সংখ্যারেখা ধরে চলি। প্রতিবার $3$ করে যোগ করি।
কতবার যোগ করার পর প্রথম ধনাত্মক সংখ্যায় পৌঁছাব? এবার একই প্রশ্ন, কিন্তু
প্রতিবার $3$ করে গুণ করলে কী হয় তা আলাদা করে বলো।

\begin{sol}
\textbf{যোগ।} পরপর পাই
\[
  -4, \quad -1, \quad 2, \quad 5, \quad \dots
\]
$-4 + 3 = -1$ (এক বার), $-1 + 3 = 2$ (দুই বার)। তাই দুইবার যোগ করার পর প্রথম
ধনাত্মক সংখ্যা $2$ পাওয়া যায়।

\medskip
\textbf{গুণ।} পরপর পাই
\[
  -4, \quad -12, \quad -36, \quad -108, \quad \dots
\]
ধনাত্মক সংখ্যা কখনোই আসবে না। ঋণাত্মক সংখ্যাকে ধনাত্মক $3$ দিয়ে গুণ করলে
চিহ্ন আলাদা, তাই ফল সবসময়ই ঋণাত্মক। সংখ্যাটা শুধু বাঁ দিকে সরতেই থাকে।

\medskip
পার্থক্যটা মনে রাখার মতো : যোগ সংখ্যারেখায় \emph{সরায়}, তাই শূন্য পেরিয়ে
যাওয়া যায়। ধনাত্মক সংখ্যা দিয়ে গুণ শুধু $0$ থেকে দূরত্ব বদলায়, দিক বদলায় না।
\end{sol}

\end{document}
